\section{One real obstruction at several joining nodes}\label{p14-sec:joins}

Fix the root $s$ on $Z$ and its transposition $\tau$, and choose once
and for all an even generator $s^+$ of the other eigenspace in $V$.
Over $x_i$ write $v_i=\partial_zs(r_i)$ and $e_i=s^+(r_i)$.
Use the quotient cubic $Y=Z/\langle\tau\rangle$. The image $y_i$
of $r_i$ is an unramified point of $Y\to C$. Trivialize $N$ in a
holomorphic frame for the complex-linear part of $D$, writing
$D=\bar\partial+B(w)\,d\bar w$, with $B$ antilinear. In that
frame put $B_i=B(x_i)$ and $L_i=N_{x_i}/\mathbb Cv_i$. We shall have
$[B_i v_i]\ne0$ and $[e_i]\ne0$ in $L_i$. Let $\psi$ generate
the adjoint line on $Y$ annihilating its branch columns in
\eqref{p14-eq:lowertransverse}. In coordinate $w$ it pairs
$a\,d\bar w$ with $\operatorname{Re}\psi(w)(a)\,dA(w)$.
Define
\begin{equation}\label{p14-eq:branchweights}
 2\overline{\beta_i}[B_i v_i]=[e_i],\qquad
 H_i(q)=\tfrac{\pi}{2}\operatorname{Re}\psi(y_i)(qv_i),
 \qquad w_i=H_i(\beta_i).
\end{equation}
Here the $w_i$ on the right are real weights. The local domain coordinate
$z$ and the target coordinate $w$ are those used in the smoothing chart
below. The same $s,s^+,\psi$ are used at all the joining flags.

\begin{proposition}\label{p14-prop:signs}
Join the sextic to the identity over a nonempty set $I$ of branch
values. Under the regularity conditions below, if
$\sum_{i\in I} w_i\ne0$, the joined map without contracted components (the bare map) is an isolated zero in the full
finite-cover chart times the parameter interval. Its local Euler change is
\begin{equation}\label{p14-eq:localsubsetsign}
              a_I=\kappa(-1)^{|I|}
                     \operatorname{sgn}\Bigl(\sum_{i\in I} w_i\Bigr),
\end{equation}
where $\kappa\in\{1,-1\}$ is common to all subsets at the crossing.
All nonempty subsums are nonzero at a generic crossing. Every
strict sign pattern of these sums can be realized on the fixed cubic.
\end{proposition}

\subsection{The derivative along the nodal stratum}
There are $3b$ independently movable simple ramification points on
$Z$, even when their images coincide. At each join there are two more
complex parameters $p,q$. In coordinates on the two branches the
family is $z y=q$, $w=(z+p)^2+y$ after translating the branch center.
Elimination gives the algebra with basis $1,z,z^2$ and relation
\begin{equation}\label{p14-eq:cubicpolynomial}
                   z(z+p)^2-wz+q=0 .
\end{equation}
At the node the relative deformation sheaf has length three: at the
central fiber its local algebra is the cokernel of the Jacobian of
$(zy,z^2+y)$, or equivalently $\mathbb C[z]/(z^3)$ after eliminating
$y$. The other simple ramification points have length one. The
relative deformation complex is supported at these finitely many
source points and the finite-map chart is smooth of complex dimension
$3b+2j$, $j=|I|$. Its central stabilizer is $\langle\tau\rangle$.
Write $P_Z=\operatorname{coker}D_Z$. Normalization and evaluation
give the joined kernel $\mathbb Rs$ and the joined cokernel $P_I$ with
\begin{equation}\label{p14-eq:joinedsequence}
 0\longrightarrow
 (\bigoplus_{i\in I}N_{x_i})/\mathbb R(e_i)_{i\in I}
 \xrightarrow{\iota}P_I\longrightarrow P_Z\longrightarrow0 .
\end{equation}
Use the convention that the ordinary forcing of a cutoff height
with value $a$ on the sextic and zero on the identity represents
$\iota(a)$.

To identify the image in $P_Z$, split under $S_3$. The trivial
part is zero. The sign part is the odd cokernel on the discriminant
double $T=Z/A_3$ and has dimension $2b$. The standard part has
multiplicity $2b+1$, its multiplicity space being
$\operatorname{coker}D_f$. These identifications follow by stabilizing
with forms away from the branches on each quotient. A bounded
invariant solution descends and extends across the punctures, and
the Fredholm indices agree by deformation to a complex-linear
operator. The odd operator on $T$ has no smooth kernel. Augment it
by a complex plane of odd poles $\chi a v_i/z$ at each ramification
point above $x_i$, as in the proof of Proposition~\ref{p14-prop:realize}.
Its resulting index is zero. We require it to be invertible. This
is open dense with the lower graph fixed. A nonzero odd polar
solution cannot be a real combination of the two lower heights on
an entire open disk: the coefficients there would be constant and
the sign representation on the punctured normal cover has zero
intersection with the standard one. An allowed zeroth-order
perturbation can prescribe arbitrary forcing on this solution on
a smaller disk. The two odd adjoint contributions add. Unique
continuation and rank increase in a finite reduction give the
assertion.

Write $B_0$ for the real space of sextic branch velocities and
increasing time. Their derivative on $s$ embeds $B_0$ as a hyperplane
in $P_Z$. To see this, first take the common shift of the three
ramification points over each $x_i$. These shifts and time preserve
$S_3$ symmetry. On the even vector their standard component is
\eqref{p14-eq:lowertransverse}; on the odd vector the same columns span
the odd standard summand since $\operatorname{End}_{S_3}V=\mathbb R$.
The additional separate shifts of the fixed flags give on the sign
part the discriminant pole columns computed above, with nonzero
factor $1/3$ from averaging over $A_3$. Indeed a branch shift $\alpha$
produces the principal part $-\alpha v_i/(2z)$, by differentiation
at fixed regular target point. Subtract the principal part and use
removability for the bounded remainder; changing the cutoff changes
the forcing by the image of a regular section. Finally the opposite
shifts at the two $\tau$-exchanged flags, applied to $(1,-1,0)$,
give one third of the common shift applied to $(1,1,-2)$.
At inertia $(13)$ and $(23)$ this follows by taking the ratios
$1/3$ and $-1/3$ of the coordinate differences. The even vector has
zero derivative at the remaining flag. These columns give exactly
the cubic branch columns on the invariant quotient. The remaining
functional $\Lambda$ is thus the half-trace to the cubic paired
with $\psi$. Extend it to $P_I$ by \eqref{p14-eq:joinedsequence}.
For the full $q_i=0$ derivative, each attaching column contributes
$v_ip_i$ in matching. They are independent modulo the one diagonal
real line since $[e_i]\ne0$. Denote this injective nodal derivative,
including time, by $J$, and its variables by $h$.


\subsection{The smoothing terms}
Let $F$ be the normalized obstruction equation, $q_i$ the complex
smoothing parameters, and $Q_i=q_i\log(1/|q_i|)$, continued by zero.
In the matching fibers the leading displacement is
$2\iota((\overline{Q}_i B_i v_i)_{i\in I})$. It vanishes under the
remaining real functional $\Lambda$. On the cone where all $|q_i|$
are comparable and the nodal parameters $h$ are
$O(R_0\log(1/R_0))$, with $R_0=\max_i|q_i|$, the next term is
$\Lambda F=\sum_i H_i(q_i)+o(R_0)$. It is the residue of the
variation of the two sextic sheets on the quotient cubic.
Appendix~\ref{app:joined-estimates} proves the two estimates,
including uniformity at the faces where some $q_i$ vanish.

\subsection{Degree and orientation}
Choose $n\in P_I$ with $\Lambda(n)=1$. The boundary degree is that
of the real-linear map
\begin{equation}\label{p14-eq:linearjoin}
 (h,q)\longmapsto Jh+2\iota\big((\bar q_iB_i v_i)_{i\in I}\big)
                                      +n\sum_i H_i(q_i).
\end{equation}
Indeed first homotope by linear interpolation to this model with a
common factor $\log(1/R_0)$ on the matching term. Replacing the separate
logarithms there by the common one costs $O(R_0)$, uniformly up to
the coordinate faces. At a sequence of noncentral zeros
\eqref{p14-eq:fullestimate} and injectivity of $J$ would force
$|h|=O(R_0\log(1/R_0))$, $R_0>0$. Projecting away the nodal
columns, then dividing by $R_0\log(1/R_0)$, puts each limiting
tuple $(q_i/R_0)$ on the real line spanned by $(\beta_i)$. Indeed
the remaining matching space is $(\bigoplus L_i)/\mathbb R([e_i])$.
Every entry of a nonzero tuple on that line is nonzero. Hence the
comparable cone applies, and division of \eqref{p14-eq:rowestimate}
by $R_0$ contradicts $\sum_i w_i\ne0$. On the model the common
positive logarithmic factor can then be changed to one. This
establishes both isolation and the asserted boundary degree.

Restore the even height as a real variable, projecting to its
evaluation vector $E_I=(e_i)$ before taking the quotient in
\eqref{p14-eq:joinedsequence}. Eliminate $B_0$ and order the remaining
rows as the matching fibers followed by the complementary normalized
line. The remaining matrix is
\[
 \begin{pmatrix}M_I&E_I\\ H&0\end{pmatrix},\qquad
 M_i(p_i,q_i)=v_i p_i+2\bar q_i B_i v_i,\qquad
 H(p,q)=\sum_i H_i(q_i).
\]
Each $M_i$ has real sign $-1$ by complex linearity in $p_i$ and
antilinearity in $q_i$ modulo $\mathbb Cv_i$. The Schur complement is
$-H M_I^{-1}E_I=-\sum_iw_i$. The sign from cancellation of the
branch-and-time block, the identity operator, and the even height is
common to all $I$: the normalization determinant is the determinant
on the two smooth components times the dual determinant of the
complex matching fibers. Every added pair $(p_i,q_i)$ and matching
fiber has real dimension four. Moving these past a fixed real line
or past each other preserves orientation. The projective Euler link
has Euler number twice the oriented degree, since its cokernel
quotient is the pullback of the tangent bundle of an even-dimensional
sphere under the normalized column. Dividing by the order-two
stabilizer leaves the signed unit \eqref{p14-eq:localsubsetsign}.

\subsection{Regularity and prescription of the weights}
The marked cubic height difference and the derivative test in
Section~\ref{p14-sec:kernels} give generically $[e_i]\ne0$. Adding the
value of $B_i$ to the marked test is submersive: after forcing off
the marks the adjoint vanishes there, the local section values and
holomorphic jets kill the marked multipliers, and the zeroth-order
coefficient can then vary freely at the center. The failure of
$[B_iv_i]\ne0$ has real codimension two.

With these jets fixed, adjoin to the cubic branch-augmented operator
one real column whose residues at the unramified $y_i$, $i\in I$,
are $\beta_iv_i$. A cutoff representative with residue $R$ there is
$R/w-B_iR\log|w|^2$. The off-pole forcing is locally
$O(1+|\log |w||)$ and its pairing with the remaining cubic adjoint
is $-2\sum_iH_i(\beta_i)$. Quotient out the known smooth kernel.
The resulting index-zero operator is invertible exactly when the sum
is nonzero. At a defect a polar solution has nonzero coefficient on
the new column. Its three conjugates at any chosen $x_i$, pulled
back to $Z$, have an order-two pole at three different flags, while
the other conjugates there have at most order one. Together with
the bounded $s_1,s_2$ these give five independent solution germs.
The constrained Petri argument of Lemma~\ref{p14-lem:kernelplane}
therefore gives a rank-increasing perturbation at a defect. More
precisely, first insert a good Petri jet on the test disk, then
apply the annihilator argument there to the three polar columns and
three local adjoint components. This gives open density with fixed
graph; intersect over the finitely many subsets. Open conditions on
a continued cubic crossing transfer to generic paths by a time
cutoff. For the countably many kernel exclusions of
Lemma~\ref{p14-lem:kernelplane} the same argument is applied to the
individual compact configuration blocks before taking the intersection.

Finally we explain how to prescribe the chamber on the fixed cover in
Proposition~\ref{p14-prop:realize}. The branch model used there starts
with $B_i=0$. Insert any bounded antilinear values with
$[B_iv_i]\ne0$ by cutoffs on disks of radius $\rho\to0$. The
perturbation tends to zero as an $L^p$ multiplier on the fixed
covers for each finite $p>2$. To restore the equation on the same
$Y$ add $2b+1$ smooth real operator variations on the regular
frame disk. Their values on $\xi$ may be chosen to span the
cokernel of $D_f$, by trace-free forcing and adjoint continuation.
The implicit function theorem with a thick real normalization now
gives coefficients tending to zero and a new normalized solution
converging to $\xi$ in $W^{1,p}$. At ramification the pulled-back
cutoff coefficient has the form $2\bar z\chi(z^2/\rho)B(z^2)$.
It is uniformly Lipschitz and $O(\sqrt\rho)$, hence small in
$C^\alpha$, $0<\alpha<1$. Interior estimates give $C^1$ convergence
of the heights at the flags. In particular the even value stays off
the complex derivative line, the odd height is $v_i z+O(|z|^{1+\alpha})$
uniformly after choosing $\alpha<1-2/p$, and regular-sheet separation
persists up to each ramification point. The branch pole columns
converge in $L^p$ since the factor $\bar z$ cancels the conjugate
polar denominator. Consequently \eqref{p14-eq:lowertransverse} remains
invertible, and its branch-annihilating adjoint, normalized by time,
converges uniformly on the cubic by the adjoint estimates. Apply the
preceding regular-disk tests to ensure the singleton $H_i(\beta_i)$
are nonzero and the discriminant augmentation is invertible.

Given real $c_i$ with all nonempty subsums nonzero, now choose
$\gamma_i$ with $H_i(\gamma_i)=c_i$. Insert bounded coefficients
$\widetilde B_i$ satisfying
$2\overline{\gamma_i}[\widetilde B_i v_i]=[e_i]$, and restore the
fixed-cover equation in exactly the same way. The normalized cubic
and discriminant pole columns and the time column converge on their
fixed sources. Thus the remaining orientation choices transport
continuously and $\beta_i\to\gamma_i$, $w_i\to c_i$. For sufficiently
small supports all desired strict inequalities hold. Restoring the
residual kernel exclusions on the regular disk preserves them.
This proves both propositions, including realizability on a
prescribed $f$. For every even $b>0$ such cubics exist over the same
$C$: send the $b$ peripheral loops to $(12)$ and one handle
generator to $(123)$, all other handle generators to the identity.
The surface relation holds and compactifying gives the required
connected holomorphic cover.
